\scanpage{87}
\mspage{4}
% archive sheet number: 86
\[
{\unsure{\Phi|\varphi}}^*_\nu\Longleftarrow E_2^{pq}=H^p(X_0,\underline\Phi^q_\nu)
\tag*{(1.14)}
\]
où $\underline\Phi^q_\nu$ est un \uline{faisceau} sur $X_0$ qui, via identification de $X_{\alpha0}$ avec $X_0$, peut s'écrire
\[
\underline\Phi^q_\nu=\varinjlim_\alpha\underline H^*_{X_{\alpha0}}(\Lambda_{\nu X_\alpha})
\tag*{(1.15)}
\]
Le calcul de ces faisceaux sur $X_0$ est \uline{de nature locale sur} $X$. [Sous réserve \deleted{que} de vérifications de passages à la limite, on conclut de \unsure{même}
\[
\Phi^*\Longleftarrow E_2^{pq}=H^p(X_0,\underline\Phi^q)
\tag*{(1.14 bis)}
\]
avec
\[
\underline\Phi^q=\varprojlim_\nu\underline\Phi^q_\nu,
\tag*{(1.15 bis)}
\]
\inserted{mais} que nous n'utiliserons pas tels quels.\deleted{]}\inserted{Il faudrait en effet \unsure{supposer}, \unsure{voire justifier ce point de vue}, que les $\underline\Phi^q_\nu$ sont \unsure{constructibles|constructible} et que le syst. projectif $(\underline\Phi^q_\nu)_{\nu\geq1}$ définit un faisceau $\ell$-adique constructible sur $X$\unsure{\dots}).}

1.16. On trouve ainsi un principe de localisation fort utile pour comparer la cohomologie de $X_s$ et de $X_{\bar{\unsure{\eta|y}}}$. Notons par exemple, \deleted{\unsure{supp}}\deleted{que le théo} (Supposant que $\ell\ne p$) que \deleted{\unsure{avec}}\inserted{si $X\longrightarrow S$ est \uline{lisse}} le théorème de pureté pour l'inclusion $s\longrightarrow S$, et le th. de changement de base lisse, nous permettent de calculer $\underline H^{\deleted{\unsure{*}}i}_{X_{\alpha0}}(\Lambda_{\nu X_\alpha})$ aux points où $X_\alpha$ est \uline{lisse} sur $S_\alpha$ : c'est nul pour $i\ne0$, et \deleted{\unsure{\simeq}} $\simeq(\mu_{\ell^\nu}^{\otimes(-1)}{}_{X_{\alpha0}})$ \deleted{\unsure{autrement}}\inserted{si $i=2$}.

On en conclut que $\underline\Phi^i_\nu=0$ si $i\ne2$, et $\underline\Phi^2_\nu=\deleted{\unsure{H}}\varinjlim_\alpha(\mu_{\ell^\nu})_{X_{\alpha0}}$. Il faut déterminer les
